\abstract{\addtocontents{toc}{\vspace{1em}} 

Predicting how rapidly stock prices will swing tomorrow is one of the trickiest challenges that financial practitioners and academic researchers face, especially in emerging economies where policy shifts and global capital flows frequently disrupt local markets. This project tackles the problem head-on for the Indian stock market by building a forecasting framework that we call the Regime-Aware Temporal Fusion Transformer (RA-TFT). The core idea behind this framework is straightforward: different market conditions---calm stretches, moderately choppy periods, sustained trending phases, and full-blown crises---require different modelling responses, so we explicitly detect these states first and then let the deep-learning model condition its forecasts on whatever state the market happens to be in.

We gathered roughly 2.65 million five-minute candlestick observations covering fourteen blue-chip stocks listed on the National Stock Exchange (NSE) of India over a span of nearly ten years. From this raw price and volume data we engineered a set of 47 features that draw on multiple well-established volatility measures: realized volatility at several time scales, lagged values motivated by the Heterogeneous Autoregressive model, jump statistics derived from bipower variation, range-based variance estimators according to Parkinson, Garman-Klass and Rogers-Satchell, GARCH(1,1) conditional variance, cross-asset correlations inspired by the Dynamic Conditional Correlation framework, and Hidden Markov Model regime labels. All these features feed into the Temporal Fusion Transformer architecture, which offers built-in interpretability through its Variable Selection Network and multi-head attention mechanism.

A carefully designed ablation study separates autoregressive persistence from genuine feature-driven prediction. Model~A, which receives the lagged target as an input, achieves a coefficient of determination $R^2$ of 0.976 with a mean absolute percentage error (MAPE) of 1.90 percent. These figures closely match the Heterogeneous Autoregressive model of Realized Volatility (HAR-RV) benchmark ($R^2 = 0.983$), confirming that the TFT can learn the autoregressive structure of volatility as effectively as the purpose-built classical model while additionally offering multi-stock generalization and built-in interpretability. Model~B, which deliberately excludes the lagged target, yields $R^2 = 0.258$, establishing that the forty-seven engineered features carry modest but genuine predictive information beyond simple persistence. Regime-conditioned evaluation reveals strong forecasting accuracy in low, normal, and high volatility states ($R^2 > 0.98$) but noticeable degradation during extreme events ($R^2 = 0.765$), highlighting the difficulty of predicting volatility during market dislocations. The Variable Selection Network independently rediscovers that lagged daily realized volatility dominates all other features by roughly two orders of magnitude, effectively learning the HAR-RV structure without being told about it.
\\
}

\textbf{\textit{Keywords}:} Volatility Forecasting, Temporal Fusion Transformer, Hidden Markov Model, GARCH, Indian Stock Market, Deep Learning, Regime Detection, NSE

\clearpage % Start a new page%